\nonumchap{Summary}
\label{ch:summary}

\draftnote{this is a placeholder skeleton --- rewrite once \cref{ch:discussion}
is filled in with concrete numbers, since the claims below should restate
quantified results, not generic statements.}

This thesis set out to determine whether sub-wavelength displacement and
material change can be recovered from time-lapse GPR surveys once classical
amplitude-based migration imaging has stopped resolving them --- the
unifying hypothesis stated in \cref{ch:introduction}.
\Cref{ch:resolution} quantified that amplitude floor for two stationary
point scatterers across three independent migration algorithms.
\Cref{ch:tl-horizontal,ch:tl-vertical} showed that the same floor limits
amplitude \emph{differencing} of a single displaced scatterer, both
laterally and vertically, with and without additive noise. A 2D Fourier
phase-plane shift estimator was then derived (\cref{ch:theory}) and
validated (\cref{ch:phase-plane}) directly on the same datasets, and applied
to a more field-realistic, spatially distributed fluid-front scenario
(\cref{ch:fluidflow}), where the same fit was shown to decouple a purely
geometric displacement from a purely dielectric material change using a
single weighted least-squares plane fit.

\draftnote{close with a final conclusion sentence stating, in one sentence,
whether the unifying hypothesis is supported, once \cref{ch:discussion} is
complete.}
